\histitem{Partial operators and the spectral theorem}

General Hilbertian spectral theory was created in large part to address the problem of the mathematical foundations of quantum mechanics, notably under the impetus of von Neumann.

We refer to the work of M. Jammer {[}35{]} for a detailed description of the development of quantum theory up to the fundamental paper by W. Heisenberg that introduced quantum mechanics {[}28{]}. Less than five years elapsed between its publication in the summer of 1924 and that of the article {[}50{]} in 1929, where von Neumann presents in a perfectly lucid and rigorous manner all the fundamental results of the theory of partial (often called "unbounded") Hermitian operators on Hilbert spaces.

By giving the first axiomatic presentation of Hilbert spaces {[}49{]}, von Neumann crowned an idea long in gestation (cf. ÉHM, p. 267), and illuminated various isomorphisms between spaces of sequences or functions well known to Schmidt, Fréchet, Fischer, and Riesz before 1910. This rise in abstraction allowed von Neumann to make considerable conceptual advances.

Such is the case with the idea of a partial operator, which seems to hover, still formless, over several of the works that, after Hilbert, took up Fredholm equations again (see above, no. 1). Hilbert had noticed that many differential equations of Sturm--Liouville type, in particular in the case of an unbounded interval, could be reduced to Fredholm equations; but to apply the integral methods, it was necessary to weaken the conditions imposed on the kernel and to leave the framework in which the results of Fredholm, Hilbert, and Schmidt could be applied. Over the years, more and more singular kernels were studied, until equations were considered for which the left-hand side makes sense only if the unknown belongs to a subspace of L²(I). In this direction, one must mention works of E. Hilb {[}30{]} in 1908, of H. Weyl on Sturm--Liouville theory {[}82{]} in 1910, and above all of T. Carleman {[}11{]} in 1923, in a very general framework.

By unifying these problems, von Neumann cast entirely new light on these classical works. He emphasized in particular the fundamental distinction between symmetric and self-adjoint operators, and interpreted the self-adjoint extensions of a symmetric operator in terms of abstract "boundary conditions," generalizing the well-known boundary conditions that appeared, for example, in Dirichlet and Neumann problems for the Laplace operator on a bounded open set of \(R^{n}\) .

At the same time, von Neumann {[}51{]} brought to the fore the notion of normal operator \(^{(5)}\) as the natural framework for spectral theory, insisting on the role of the commutative algebra generated by the operator and its adjoint.

Relying on this precise mathematical formalism, which also revealed the equivalence between the viewpoints of Heisenberg and Schrödinger, and on the Copenhagen interpretation of experimental procedures and results, non-relativistic quantum mechanics from then on reached a state of refinement that has hardly evolved since; the validity of this physical theory, despite its surprising, even paradoxical-looking consequences, has since been constantly confirmed by experiment, and with remarkable precision.

From a formal point of view, the foundations of spectral theory were laid by von Neumann. In parallel, after the presentation of the latter's first ideas, M. Stone had pursued similar research between 1928 and 1930 and obtained closely related results. The clear and complete presentation of the known results that Stone published in 1932 {[}72{]} had an important influence on the diffusion of Hilbertian spectral theory.

Nevertheless, other improvements in the presentation have been important for the diffusion of these results. Indeed, the spectral theorem (cf. th. 1 of IV, p. 266), as presented by von Neumann, was long considered very difficult, largely because it was stated in the framework of vector-valued measures.

Later, P. Halmos {[}27{]}, by emphasizing the interpretation of the spectral theorem by means of multiplication operators on \(L^{2}\) spaces, revealed its essentially elementary aspect in the case of bounded operators. Quite recently, S. Woronowicz (in the somewhat different framework of regular operators on Hilbert modules over C*-algebras {[}92{]}) brought a simplification by introducing

the "bornification" (cf. no. 2 of IV, p. 265), which allows one to treat normal partial operators as simply as self-adjoint operators.

Moreover, if D is a formally symmetric scalar differential operator on an open set U of \(R^{n}\) , the study of the self-adjoint extensions of D to \(L^{2}(U)\) is naturally linked to the study of distributions on U (in the case of the Laplacian, cf. no. 6 of IV, p. 242). We refer to the work of J. Dieudonné {[}16, ch. VII{]} for the history of the latter --- a sinuous history that unravels shortly after the period we have just evoked, with decisive contributions of S. Sobolev {[}69{]} in 1936, and of L. Schwartz {[}65{]} ten years later.

From a purely mathematical point of view, spectral theory interacts in a particularly fruitful way with Riemannian geometry. Thus, taking up a question of S. Bochner, M. Kac gave in 1966 a strikingly vivid presentation {[}36{]} of the problem of determining which geometric properties of a Riemannian manifold can be deduced from the spectral properties of the Laplace operator. This question had been anticipated by the physicist A. Schuster {[}64{]} in 1882: To find out the different tunes sent out by a vibrating system is a problem which may or may not be solvable in certain special cases, but it would baffle the most skillful mathematician to solve the inverse problem and to find out the shape of a bell by means of the sounds which it is capable of giving out. (“Determining the harmonics emitted by a vibrating system is a problem which may or may not be solvable in certain special cases, but the most skillful mathematician would be baffled if he had to solve the inverse problem and determine the shape of a bell from the sounds it can emit”).

